
\section{Frequently Asked Questions}

\begin{quote}
\textit{\textbf{Audience note:} This section is written in plain English for
non-technical stakeholders, business contacts, and junior developers who need practical
clarity without deep protocol jargon. It deliberately departs from the RFC tone used in
all other sections of this specification. If a question involves a normative requirement,
the relevant section of the spec is linked for reference.}
\end{quote}

\paragraph{Q1: Why doesn't Canvasflow support immediate token revocation?}

\textbf{Short answer:} It's a deliberate design choice, not a missing feature.

Here's the thing: when your user logs in, your IdP (your identity system) creates a
signed token and hands it to our app. That token is like a signed ticket --- once it's
been signed and issued, Canvasflow has no way to ``un-sign'' it. We're only reading the
ticket, not the box office.

To immediately revoke access, you would need to call Canvasflow's server on every single
API request to check ``is this token still valid?'' --- but that would mean Canvasflow has
to phone home to your system tens or hundreds of times per session. That creates a
dependency: if your system is slow or down, your users can't use Canvasflow either. We
chose not to build that dependency.

Instead, tokens are short-lived --- they expire after at most one hour. If you remove a
user's access in your system today, they lose it at their next login or when their token
expires, whichever comes first. That's the revocation window.

\textbf{If you need faster revocation:} revoke the user's refresh token in your IdP
right now. That stops them from getting any new access tokens. The old access token still
works until it expires (up to one hour), but after that they're locked out and can't
refresh.



See \hyperref[sec:out-of-scope]{\S1.4} and \hyperref[sec:freshly-evaluated-entitlements]{\S8.3} for the technical details.

\paragraph{Q2: What does an empty or missing entitlements claim mean? Will my users see
an error?}

\textbf{Short answer:} No error. They just see the free content, not the premium content.

If a user's token has no entitlements --- because the claim is empty, it's missing
entirely, or all the values are in the wrong format --- Canvasflow treats them as a
free-tier user. They can log in just fine. They'll see everything Canvasflow marks as
free content. They just won't see anything behind a paywall.

A few things worth knowing:
\begin{itemize}
  \item \textbf{Empty array (\texttt{[]})} is valid and intentional. If someone is on a
    free plan, have your hook emit an empty array. That's the correct signal.
  \item \textbf{Missing claim entirely} works the same as empty, but we log it on our
    side as a possible misconfiguration --- because we can't tell if it's intentional or a
    bug in your hook.
  \item \textbf{Wrong format} (more on this in Q9) also results in free content, not an
    error.
\end{itemize}

See \hyperref[sec:free-content-baseline]{\S4.4} for the full details.

\pagebreak

\paragraph{Q3: Which claim name should I use for entitlements?}

\textbf{Short answer:} Use \texttt{cf:entitlements}. It works everywhere.

There are three claim names Canvasflow accepts:
\begin{enumerate}
  \item \texttt{cf:entitlements} --- \textbf{use this one}. Works on every IdP, including
    Auth0.
  \item \texttt{resources} --- a legacy claim name, still supported for backward
    compatibility. If you're starting fresh, skip it.
  \item \texttt{entitlements} --- the official standards-based name (from an internet
    standard called RFC~9068). Works on most IdPs, but Auth0 blocks it. Risky choice.
\end{enumerate}

If you're using Auth0, you must use \texttt{cf:entitlements} --- the other names will be
quietly dropped by Auth0 and you won't get an error, your claim just won't show up.

If all three are somehow present in a single token (unusual, but possible if you have
multiple hooks), \texttt{cf:entitlements} wins and the others are ignored.

See \hyperref[sec:claim-name-hierarchy]{\S4.3} for the precise rules.

\paragraph{Q4: Why is my Auth0 token opaque? Why is my custom claim missing?}

These are the two most common Auth0 ``gotchas.'' Here's the checklist:

\textbf{Token is opaque (looks like a random string, not a JSON Web Token):}
\begin{itemize}
  \item You need to register an \textbf{API} in Auth0, not just an Application. The
    API's Identifier must be set to \texttt{canvasflow-api}.
  \item Your authorization request from the SPA must include
    \texttt{audience=canvasflow-api}.
  \item Without both of these, Auth0 gives you a short opaque token instead of a JWT.
    Canvasflow can't read opaque tokens.
\end{itemize}

\textbf{Custom claim is missing from the token:}
\begin{itemize}
  \item You MUST use \texttt{cf:entitlements} as the claim name. Auth0 silently drops
    claim names like \texttt{entitlements} or \texttt{resources} from access tokens
    (they're on Auth0's blocklist for non-namespaced claims).
  \item Make sure your Action is deployed \textbf{and} attached to the post-login
    Trigger in \textbf{Actions $\to$ Triggers}. A deployed-but-unattached Action does
    nothing.
  \item Don't rely on \texttt{console.log} to debug --- it doesn't show up in Auth0's
    monitoring logs. Decode the actual issued token instead.
\end{itemize}

See \hyperref[sec:auth0-notes]{\S6.3} for the full Auth0 operational notes.

\paragraph{Q5: Why should I use JWKS instead of a shared secret?}

\textbf{Short answer:} With JWKS, a Canvasflow breach can't be used to fake your users'
tokens.

With a shared secret (HS256), both you and Canvasflow hold the same key. If Canvasflow's
systems were ever breached, an attacker who got that secret could forge tokens that look
completely legitimate --- giving themselves access to your users' content.

With JWKS (RS256 or ES256), you keep your private key locked inside your IdP. It never
leaves. Canvasflow only ever gets your public key, which can verify signatures but can't
create them. A Canvasflow breach gives an attacker your public key --- which is already
public --- and nothing useful.

The other benefits:
\begin{itemize}
  \item \textbf{Key rotation is painless.} Post a new key to your JWKS URL, keep the old
    one for an hour, done. No calls to us, no downtime.
  \item \textbf{No secret to manage on our end.} Every shared secret we hold is a
    security liability (storage, access controls, breach notifications). JWKS eliminates
    that.
  \item \textbf{Your IdP probably already does this.} Cognito, Okta, Azure AD, and Auth0
    all publish a JWKS endpoint by default.
\end{itemize}

See \hyperref[sec:signing-path-requirements]{\S5.2} for the full technical comparison.

\paragraph{Q6: When does a user gain or lose access after I change their entitlements?}

\textbf{At their next login or token refresh --- not immediately.}

When a user logs in or their session refreshes, your hook runs again, reads their current
entitlements from your system, and puts the result in a fresh token. So:
\begin{itemize}
  \item \textbf{Add an entitlement:} the user gets access on their next login or refresh.
    If they're already logged in, they might wait up to one hour (the token lifetime).
  \item \textbf{Remove an entitlement:} same deal. The change takes effect on their next
    login or refresh.
  \item \textbf{Revoke now:} invalidate their refresh token in your IdP. They can't get
    a new access token. Their current access token still works until it expires (up to
    one hour), then they're fully locked out.
\end{itemize}


This is a known trade-off of this architecture. See 
\hyperref[sec:freshly-evaluated-entitlements]{\S8.3} and 
\hyperref[sec:session-lifecycle]{\S3.2.6} for more.

\paragraph{Q7: What happens if I accidentally put the entitlements claim in the ID token
instead of the access token?}

\textbf{Your users can log in but won't see any premium content.}

The access token is what Canvasflow reads. The ID token is for your app's UI (displaying
the user's name, picture, etc.) and is never sent to our API. If the claim ends up in
the ID token, Canvasflow receives an access token with no entitlements --- which is treated
as a free-tier user.

Your users won't get an error. They'll just wonder why the premium content isn't showing
up, even though they're logged in.

The Canvasflow OAuth Validator (\url{https://oauth.canvasflow.work}) catches this exact
mistake before go-live. If you run the validator and V3 shows an empty entitlement set
when you expected content, check whether the claim ended up in the wrong token.



This is the most common hook misconfiguration --- see \hyperref[sec:hook-rules]{\S6.4}, Rule~3.

\paragraph{Q8: My users authenticate but only see free content. What's wrong?}

This usually means Canvasflow isn't seeing any valid entitlement values in the access
token. Work through this checklist:

\begin{enumerate}
  \item \textbf{Is the claim in the access token?} Decode the access token (not the ID
    token) at \url{https://jwt.io}. Look for \texttt{cf:entitlements},
    \texttt{resources}, or \texttt{entitlements}. If it's not there, your hook isn't
    running or is targeting the wrong token (see Q7 and \hyperref[sec:hook-rules]{\S6.4}).

  \item \textbf{Are the values in \texttt{cf:<type>:<id>} format?} If you see values
    like \texttt{issue:123} or \texttt{456} --- without the \texttt{cf:} prefix --- those
    are silently ignored. Only values starting with \texttt{cf:} count. See Q9 below.

  \item \textbf{Is the claim name allowed on your IdP?} If you're on Auth0 and using
    \texttt{entitlements} as the claim name, Auth0 silently drops it. Switch to
    \texttt{cf:entitlements} (see Q3).

  \item \textbf{Did you run the validator?} The Canvasflow OAuth Validator
    (\url{https://oauth.canvasflow.work}) flags all of these problems with specific
    explanations. Run it --- it's the fastest way to diagnose this.
\end{enumerate}

\paragraph{Q9: My claim has values in it but some content still isn't accessible. Why?}

\textbf{The values probably aren't in the right format.}

Canvasflow only pays attention to entitlement values that start with \texttt{cf:} and
match the format \texttt{cf:<type>:<id>}. For example:

\begin{table}[htbp]
\centering
\begin{tabular}{@{}llp{7cm}@{}}
\toprule
\textbf{Value} & \textbf{Accepted?} & \textbf{Why} \\
\midrule
\texttt{cf:issue:123} & $\checkmark$ Yes & Correct format \\
\texttt{issue:123}    & $\times$ No      & Missing \texttt{cf:} prefix \\
\texttt{123}          & $\times$ No      & Missing prefix and type \\
\texttt{product:789}  & $\times$ No      & Tenant-internal format, no \texttt{cf:} prefix \\
\bottomrule
\end{tabular}
\caption{Entitlement value format: accepted vs.\ rejected}
\end{table}

Values without the \texttt{cf:} prefix are silently ignored --- they don't cause an error
and they don't grant any access. This is intentional: we don't want one bad value in the
array to lock out the entire user.

When this happens, we log the ignored values on our side correlated with the token's
unique ID (called \texttt{jti}). If you contact Canvasflow support with the approximate
time and a user identifier, we can pull those logs and show you exactly which values were
ignored and why.

Fix: update your hook's mapping table to emit \texttt{cf:issue:123} instead of
\texttt{issue:123}. See \hyperref[sec:mandatory-value-format]{\S4.2} for the format specification.

\paragraph{Q10: What if I rotate my JWKS keys without maintaining an overlap window?}

\textbf{You'll get a brief burst of failed requests, but it self-heals.}

Here's what happens: Canvasflow caches your public keys for up to one hour. If you
remove the old key before that hour is up, new tokens (signed with the new key) arrive
at Canvasflow, which looks in its cache for the old key ID (\texttt{kid}) and can't find
the new one. Signature verification fails $\to$ 401 for the user.

However, an unknown \texttt{kid} immediately triggers a fresh fetch of your JWKS
endpoint. So the very next request after the failure will likely succeed, because
Canvasflow will re-fetch and find your new key.

Still, you should keep the old key in your JWKS for at least one hour after you start
signing with the new one. This is the overlap window we ask for, and it prevents the
brief 401 burst entirely. No-overlap rotations aren't catastrophic, but they're
preventable. See \hyperref[sec:jwks-endpoint]{\S5.2.1} for the JWKS rotation requirements.

\pagebreak

\paragraph{Q11: Can the \texttt{aud} claim be an array? I'm seeing an array with two
values.}

\textbf{Yes, that's fine. We look for \texttt{canvasflow-api} inside the array, not as
the only value.}

JWT allows \texttt{aud} to be either a single string or an array of strings. Auth0 in
particular puts both \texttt{canvasflow-api} and the Auth0 \texttt{/userinfo} URL into
the \texttt{aud} array when you request the \texttt{openid} scope. This is
spec-compliant.

Canvasflow validates that \texttt{canvasflow-api} is \textbf{present} in the \texttt{aud}
value --- it doesn't require it to be the only value. So an \texttt{aud} of
{\scriptsize\texttt{["canvasflow-api",\ "https://tenant.auth0.com/userinfo"]}} passes just fine.
See \hyperref[sec:required-payload-claims]{\S5.4} (the \texttt{aud} row in the required payload claims table).

\paragraph{Q12: I use Cognito and my token has \texttt{typ: "JWT"}, not
\texttt{typ: "at+jwt"}. Will it be rejected?}

\textbf{No --- Cognito gets a special accommodation for this.}

AWS Cognito doesn't follow the latest JWT access token standard 
(\href{https://tools.ietf.org/html/rfc9068}{RFC~9068}) in the way it
sets the \texttt{typ} header. Instead of \texttt{typ: "at+jwt"}, Cognito sets
\texttt{typ: "JWT"} and adds a separate claim called \texttt{token\_use: "access"}.

At onboarding, you declare that you're using Cognito. Canvasflow marks your tenant
configuration with a flag that tells our validator: for this tenant, accept
\texttt{typ: "JWT"} and check \texttt{token\_use: "access"} instead of the standard
\texttt{typ} check. The end result is the same level of protection against token
confusion attacks --- just using Cognito's own mechanism instead of the standard one.

So: declare Cognito at onboarding and you're covered. See 
\hyperref[sec:token-type]{\S5.3.1} for the technical details of the Cognito 
accommodation.

\paragraph{Q13: Does Canvasflow ever see my users' passwords or personal information?}

\textbf{No, never.}

Your IdP handles all authentication. When your user types their password, enters an MFA
code, or clicks ``Sign in with Google,'' that happens entirely inside your system.
Canvasflow never sees any of it.

By the time anything reaches Canvasflow, it's already a signed token --- a short document
that says ``this user is authenticated and has access to these resources.'' No password,
no credentials, no authentication details.

We also require that access tokens not include personal information like names, email
addresses, or physical addresses (\hyperref[sec:prohibited]{\S5.5}). The ID token --- which may contain user display
info --- is consumed by the SPA in your users' browsers and is never sent to Canvasflow's
servers.

